\documentclass[11pt,a4paper]{article}

\input{glyphtounicode}
\pdfgentounicode=1
\usepackage[utf8]{inputenc}
\usepackage[T1]{fontenc}
\usepackage{lmodern}
\usepackage[spanish,es-tabla]{babel}
\usepackage{geometry}
\usepackage{booktabs}
\usepackage{array}
\usepackage{longtable}
\usepackage{graphicx}
\usepackage{hyperref}
\usepackage{xcolor}

\geometry{margin=2.4cm}
\let\oldsection\section
\renewcommand{\section}{\clearpage\oldsection}
\hypersetup{
    colorlinks=true,
    linkcolor=black,
    citecolor=black,
    urlcolor=blue
}

\title{\textbf{Sistema multiagente con modelos de lenguaje para an\'alisis financiero hist\'orico}}
\author{M\'aster en Deep Learning\\Proyectos de Deep Learning}
\date{Curso 2025--2026}

\begin{document}

\maketitle

\begin{abstract}
Este documento presenta el dise\~no de un proyecto software que integra modelos de Deep Learning, en concreto modelos de lenguaje de gran escala, dentro de un flujo multiagente para an\'alisis financiero hist\'orico. El objetivo no es construir un asesor financiero ni predecir precios futuros, sino desarrollar un sistema que permita transformar consultas estructuradas sobre activos financieros en an\'alisis reproducibles, c\'odigo ejecutable y explicaciones comprensibles. El trabajo parte del MVP desarrollado para el TFM: una arquitectura modular que separa ingesta, enrutado, planificaci\'on, generaci\'on de c\'odigo, ejecuci\'on e interpretaci\'on, manteniendo una base determinista y una integraci\'on opcional de modelos LLM mediante APIs compatibles con OpenAI.
\end{abstract}

\tableofcontents
\newpage

\section{Introducci\'on}

\subsection{T\'itulo del proyecto}

El proyecto se titula \textit{Sistema multiagente con modelos de lenguaje para an\'alisis financiero hist\'orico}. La propuesta se enmarca en el desarrollo de un prototipo software que utiliza agentes especializados, ejecuci\'on program\'atica de c\'odigo y modelos de lenguaje para automatizar an\'alisis descriptivos sobre datos financieros.

\subsection{Problema que se quiere resolver}

El an\'alisis financiero hist\'orico combina tareas heterog\'eneas: localizar y cargar datos, validar entradas, seleccionar el tipo de an\'alisis, calcular m\'etricas, generar c\'odigo, ejecutar dicho c\'odigo, interpretar resultados y presentar una respuesta \'util para el usuario. En muchos sistemas basados en LLM, estas fases se concentran en una \'unica llamada generativa. Ese enfoque puede producir respuestas fluidas, pero dificulta comprobar qu\'e datos se han usado, qu\'e c\'alculos se han realizado y d\'onde se ha podido introducir un error.

El problema que se aborda es la falta de trazabilidad y control en prototipos de an\'alisis financiero asistidos por inteligencia artificial. El sistema separa el razonamiento textual de los c\'alculos verificables, de modo que el modelo de lenguaje pueda aportar valor en la planificaci\'on o la explicaci\'on sin sustituir las validaciones ni la ejecuci\'on controlada.

\subsection{Contexto y motivaci\'on}

Los modelos de lenguaje actuales han demostrado capacidad para razonar sobre instrucciones, generar c\'odigo y redactar explicaciones adaptadas al usuario. Sin embargo, en dominios sensibles como el financiero resulta peligroso aceptar una respuesta textual sin evidencia intermedia. Incluso cuando el objetivo es descriptivo, una interpretaci\'on err\'onea de rentabilidades, volatilidad o periodos temporales puede inducir decisiones equivocadas.

La motivaci\'on del proyecto es construir una arquitectura en la que los LLM se integren dentro de un flujo gobernado por reglas, contratos de datos y ejecuci\'on reproducible. La versi\'on inicial del MVP funciona de manera determinista; despu\'es se activa el modelo de lenguaje por fases, empezando por la interpretaci\'on final. As\'i se puede medir el impacto del componente generativo sin comprometer la estabilidad del sistema.

\subsection{Objetivos del proyecto}

El objetivo general es desarrollar y validar un prototipo multiagente capaz de ejecutar an\'alisis financieros hist\'oricos a partir de entradas estructuradas y datos en CSV, incorporando modelos de lenguaje de forma gradual y controlada.

Los objetivos espec\'ificos son:

\begin{itemize}
    \item Dise\~nar un workflow modular que separe ingesta, enrutado, planificaci\'on, generaci\'on de c\'odigo, ejecuci\'on e interpretaci\'on.
    \item Implementar familias de agentes especializadas para distintas intenciones anal\'iticas.
    \item Mantener un modo determinista reproducible que sirva como l\'inea base.
    \item Integrar modelos de lenguaje mediante APIs compatibles con OpenAI, con activaci\'on granular y fallback.
    \item Evaluar el sistema con casos funcionales y escenarios de error controlado.
    \item Delimitar el alcance \'etico del sistema para evitar recomendaciones de inversi\'on o predicciones no justificadas.
\end{itemize}

\section{Estado del arte}

\subsection{Modelos de lenguaje y uso de herramientas}

La investigaci\'on reciente sobre LLM ha mostrado que estos modelos pueden combinar razonamiento textual con acciones externas. ReAct propone intercalar razonamiento y acciones para que el modelo pueda planificar, consultar herramientas y actualizar su respuesta a partir de observaciones \cite{react}. Toolformer explora c\'omo un modelo puede aprender a utilizar APIs externas para complementar su generaci\'on con resultados procedentes de herramientas \cite{toolformer}. Estas ideas son relevantes para el proyecto porque el an\'alisis financiero exige acceder a datos y ejecutar c\'alculos, no solo producir texto.

El TFM adopta una visi\'on prudente de este paradigma. El LLM no act\'ua como una caja negra que decide todo el flujo, sino como un componente dentro de una arquitectura con fases separadas. La planificaci\'on, la generaci\'on de c\'odigo y la interpretaci\'on pueden apoyarse en modelos de lenguaje, pero los datos, las validaciones y los artefactos de ejecuci\'on permanecen observables.

En la misma l\'inea, los sistemas multiagente recientes dividen una tarea compleja en agentes conversacionales o funcionales con responsabilidades acotadas. AutoGen, por ejemplo, propone aplicaciones LLM compuestas por varios agentes que colaboran mediante mensajes y herramientas \cite{autogen}. Esta idea encaja con el proyecto, aunque aqu\'i se aplica con m\'as control: los agentes no son procesos aut\'onomos sin restricciones, sino m\'odulos especializados con contratos de entrada y salida.

\subsection{LLM en el dominio financiero}

El uso de modelos de lenguaje en finanzas se ha acelerado con trabajos como FinGPT, que plantea modelos abiertos adaptados al sector financiero \cite{fingpt}. Estos enfoques buscan aprovechar fuentes financieras, noticias, datos de mercado y tareas de an\'alisis espec\'ificas del dominio. No obstante, el dominio financiero introduce riesgos particulares: interpretaciones con apariencia de consejo profesional, sesgos en datos hist\'oricos, cambios de r\'egimen de mercado y alta sensibilidad a errores num\'ericos.

El proyecto se diferencia de un sistema de predicci\'on de mercado. Su objetivo es descriptivo y educativo: calcular m\'etricas hist\'oricas, comparar activos, estimar retornos y riesgo, y generar explicaciones sobre datos pasados. Esta delimitaci\'on reduce el riesgo de uso indebido y permite centrar el trabajo en la arquitectura software. Frente a modelos financieros orientados a recomendaci\'on o predicci\'on, el criterio de calidad principal no es anticipar precios, sino mantener fidelidad a los datos, trazabilidad de los c\'alculos y ausencia de recomendaciones de inversi\'on.

\subsection{MLOps y trazabilidad}

Las pr\'acticas MLOps proponen registrar experimentos, datos, par\'ametros, artefactos y m\'etricas para mejorar la reproducibilidad del ciclo de vida de machine learning. Plataformas como MLflow se han popularizado para gestionar ejecuciones, modelos y resultados \cite{mlflow}. En el proyecto, esta filosof\'ia se traslada a un MVP ligero: cada ejecuci\'on genera salidas estructuradas, scripts y logs que permiten revisar la relaci\'on entre entrada, c\'alculo y respuesta.

\section{Datos}

\subsection{Fuente de los datos}

Los datos utilizados en el MVP proceden de series financieras hist\'oricas descargadas y almacenadas localmente en CSV. El MVP desarrollado utiliza ejemplos de activos como NVIDIA, AMD, Apple, Bitcoin, QQQ, SPY, oro, EUR/USD y S\&P 500. La fuente prevista para ampliar el sistema es \texttt{yfinance}, junto con otros proveedores p\'ublicos o APIs financieras cuando sea necesario contrastar datos.

\subsection{Formato de los datos}

El formato principal es CSV tabular. Cada fichero representa una o varias series temporales con columnas de precio y volumen, fechas o marcas temporales, y metadatos asociados en ficheros JSON. La entrada normalizada del sistema incluye:

\begin{itemize}
    \item consulta original;
    \item intenci\'on anal\'itica;
    \item lista de tickers;
    \item rango temporal o periodo;
    \item intervalo de muestreo;
    \item rutas a los CSV que se deben analizar.
\end{itemize}

\subsection{Volumen esperado}

La versi\'on actual incluye ocho CSV de datos brutos, con aproximadamente 4.345 filas y un tama\~no total cercano a 540 KB. Es un volumen peque\~no, adecuado para validar el MVP sin infraestructura pesada. En una versi\'on ampliada, el sistema podr\'ia trabajar con varios a\~nos de datos diarios o intrad\'ia, aumentando a cientos de miles o millones de filas si se incorporan m\'as activos, intervalos horarios y ventanas temporales largas.

\begin{table}[h]
\centering
\begin{tabular}{lrr}
\toprule
\textbf{Conjunto} & \textbf{Filas aprox.} & \textbf{Tama\~no aprox.} \\
\midrule
Apple tres meses & 62 & 6,5 KB \\
NVIDIA vs AMD dos a\~nos & 502 & 102,6 KB \\
NVIDIA cinco a\~nos & 1.257 & 138,3 KB \\
Bitcoin 2024 & 367 & 33,0 KB \\
S\&P 500 desde 2020 & 1.560 & 157,6 KB \\
EUR/USD diez d\'ias a 1h & 226 & 26,9 KB \\
QQQ y SPY 2024 & 253 & 50,5 KB \\
Oro una semana a 1h & 118 & 11,8 KB \\
\bottomrule
\end{tabular}
\caption{Volumen de datos brutos disponible en el MVP.}
\end{table}

\subsection{Necesidad de datos adicionales}

Para un prototipo acad\'emico, los datos actuales son suficientes para validar el flujo. Para una versi\'on m\'as robusta har\'ian falta:

\begin{itemize}
    \item mayor cobertura temporal para probar diferentes ciclos de mercado;
    \item activos de distintos sectores y clases: acciones, \'indices, divisas, materias primas y criptoactivos;
    \item datos corporativos o fundamentales si se quisiera ampliar el an\'alisis m\'as all\'a de precios;
    \item datos de referencia para comparar resultados con benchmarks externos;
    \item registros de uso real para evaluar consultas frecuentes y errores de usuario.
\end{itemize}

\subsection{Preprocesamiento necesario}

El preprocesamiento incluye normalizaci\'on de nombres de columnas, conversi\'on de fechas, ordenaci\'on temporal, comprobaci\'on de valores ausentes, verificaci\'on de duplicados, validaci\'on de tickers y alineaci\'on de series cuando se comparan varios activos. Tambi\'en se requiere controlar cambios de escala, splits, dividendos y diferencias de calendario entre mercados.

\subsection{Evaluaci\'on de la calidad de los datos}

La calidad se evaluar\'a mediante reglas autom\'aticas: presencia de columnas obligatorias, continuidad temporal razonable, porcentaje de valores ausentes, coherencia de precios positivos, detecci\'on de outliers extremos y consistencia entre ficheros de datos y metadatos. Las incidencias se guardar\'an como avisos dentro del estado de ejecuci\'on para que el usuario conozca las limitaciones del an\'alisis.

\subsection{Fuentes adicionales}

Si los datos locales fueran insuficientes se podr\'ian incorporar APIs como Alpha Vantage, Stooq, Nasdaq Data Link o proveedores institucionales. Para mantener reproducibilidad, cada descarga deber\'ia registrar fecha de extracci\'on, proveedor, ticker, intervalo, zona horaria y versi\'on del proceso de limpieza.

\section{Enfoque y modelado}

\subsection{Tipos de modelos a evaluar}

El proyecto eval\'ua principalmente modelos de lenguaje de gran escala como componente de Deep Learning. No se plantea entrenar una red neuronal financiera desde cero en la fase inicial, porque el objetivo no es predecir precios sino orquestar an\'alisis reproducibles. Los modelos a evaluar son:

\begin{itemize}
    \item LLM servidos mediante APIs compatibles con OpenAI;
    \item modelos abiertos ejecutados en infraestructura universitaria mediante vLLM;
    \item modelos comerciales o externos como Groq y Gemini para comparar latencia, calidad y estabilidad;
    \item componentes deterministas como l\'inea base no neuronal.
\end{itemize}

\subsection{Justificaci\'on del enfoque elegido}

El enfoque multiagente se justifica por la necesidad de separar responsabilidades. Un \'unico prompt podr\'ia generar una respuesta, pero ser\'ia m\'as dif\'icil auditarlo. La arquitectura propuesta distingue entre:

\begin{enumerate}
    \item ingesta y validaci\'on de datos;
    \item enrutado de la intenci\'on;
    \item selecci\'on de agente especializado;
    \item generaci\'on de plan anal\'itico;
    \item generaci\'on o selecci\'on de c\'odigo;
    \item ejecuci\'on controlada;
    \item interpretaci\'on final.
\end{enumerate}

Esta separaci\'on permite activar el LLM solo donde aporte valor. Por ejemplo, en una primera fase se usa para redactar la interpretaci\'on a partir de resultados ya calculados. En fases posteriores podr\'ia participar en la planificaci\'on o en la generaci\'on de c\'odigo, siempre con validaciones y fallback.

\begin{figure}[h]
\centering
\small
\resizebox{\textwidth}{!}{%
\begin{tabular}{ccccccc}
\fbox{\parbox{2.2cm}{\centering Entrada\\estructurada}} &
$\rightarrow$ &
\fbox{\parbox{2.2cm}{\centering Ingesta y\\validaci\'on}} &
$\rightarrow$ &
\fbox{\parbox{2.2cm}{\centering Router de\\intenci\'on}} &
$\rightarrow$ &
\fbox{\parbox{2.2cm}{\centering Agente\\especialista}} \\
&&&&&& \\
\fbox{\parbox{2.2cm}{\centering Respuesta\\final}} &
$\leftarrow$ &
\fbox{\parbox{2.2cm}{\centering Interpretaci\'on\\LLM o fallback}} &
$\leftarrow$ &
\fbox{\parbox{2.2cm}{\centering Ejecuci\'on\\controlada}} &
$\leftarrow$ &
\fbox{\parbox{2.2cm}{\centering Plan y c\'odigo\\Python}} \\
\end{tabular}}
\caption{Arquitectura funcional del MVP multiagente.}
\end{figure}

\subsection{Modelos base}

Se utilizar\'an modelos preentrenados mediante transfer learning o APIs ya disponibles. No se propone entrenar desde cero por coste, volumen de datos y falta de necesidad para el objetivo del proyecto. La configuraci\'on actual contempla perfiles para Groq, Gemini y un endpoint universitario vLLM con el modelo \texttt{openai/gpt-oss-20b}. La capa LLM se activa mediante variables de entorno y puede limitarse a interpretaci\'on, planificaci\'on o generaci\'on de c\'odigo.

\subsection{Frameworks y librer\'ias}

El sistema se implementa en Python. Las librer\'ias principales son:

\begin{itemize}
    \item \texttt{pandas} para carga y transformaci\'on de datos;
    \item \texttt{numpy} para c\'alculos num\'ericos;
    \item \texttt{matplotlib} para gr\'aficos;
    \item \texttt{yfinance} para obtenci\'on de datos;
    \item clientes compatibles con OpenAI para APIs LLM;
    \item \texttt{unittest} o \texttt{pytest} para pruebas automatizadas;
    \item MLflow o Weights \& Biases como extensi\'on MLOps futura.
\end{itemize}

\section{Entrenamiento}

\subsection{Divisi\'on del dataset}

Dado que el sistema no entrena inicialmente un modelo predictivo sobre series temporales, no se aplica una divisi\'on cl\'asica de entrenamiento, validaci\'on y test para aprender par\'ametros. En su lugar se define una separaci\'on por casos de evaluaci\'on:

\begin{itemize}
    \item conjunto de desarrollo: consultas usadas para construir y depurar el flujo;
    \item conjunto de validaci\'on: consultas representativas de las seis intenciones soportadas;
    \item conjunto de test: consultas nuevas y escenarios de error para comprobar robustez.
\end{itemize}

Si en una fase posterior se entrenara o ajustara un clasificador de intenciones, s\'i se usar\'ia una divisi\'on estratificada por tipo de consulta. Para series temporales, cualquier divisi\'on predictiva deber\'ia respetar el orden temporal y evitar fugas de informaci\'on.

En este trabajo, el apartado de entrenamiento se entiende como preparaci\'on y validaci\'on del sistema de Deep Learning integrado, no como ajuste de pesos. La unidad experimental es una consulta completa: entrada estructurada, agente seleccionado, c\'odigo generado, ejecuci\'on y respuesta final. Esta decisi\'on mantiene el trabajo alineado con el objetivo de la asignatura, porque el componente neuronal existe como LLM preentrenado integrado en un proyecto software real.

\subsection{M\'etricas de evaluaci\'on}

Las m\'etricas se adaptan al tipo de componente:

\begin{itemize}
    \item \textbf{Workflow}: tasa de ejecuci\'on correcta, errores controlados y cobertura de intenciones.
    \item \textbf{C\'alculo}: coincidencia con resultados esperados, ausencia de excepciones y coherencia de unidades.
    \item \textbf{LLM}: utilidad de respuesta, fidelidad a los datos, ausencia de recomendaciones financieras y tasa de fallback.
    \item \textbf{Sistema}: latencia, coste por ejecuci\'on, trazabilidad de artefactos y estabilidad ante entradas inv\'alidas.
\end{itemize}

\subsection{Optimizaci\'on}

En la fase actual no se optimizan pesos neuronales. La optimizaci\'on se centra en el dise\~no del pipeline: reducir errores de enrutado, mejorar contratos de salida, limitar llamadas innecesarias al LLM y ajustar prompts para que el modelo respete los resultados calculados. Si se ajustara un modelo peque\~no para clasificaci\'on de intenciones, se emplear\'ian t\'ecnicas est\'andar como regularizaci\'on, early stopping y b\'usqueda de hiperpar\'ametros.

\subsection{Hardware necesario}

El MVP determinista se ejecuta en CPU. La integraci\'on LLM mediante APIs externas no requiere GPU local. Si se usa el endpoint universitario con vLLM, el coste de GPU se desplaza al servidor. Para servir un modelo abierto de forma local o privada ser\'ia recomendable una GPU con memoria suficiente, dependiendo del tama\~no del modelo y de la cuantizaci\'on.

\subsection{Uso de MLOps y seguimiento de experimentos}

El seguimiento de experimentos registrar\'a proveedor LLM, modelo, par\'ametros de activaci\'on, caso evaluado, latencia, estado, avisos, uso de fallback y respuesta final. En una versi\'on ampliada, MLflow permitir\'ia centralizar estos artefactos y comparar perfiles deterministas, Groq, Gemini y vLLM universitario.

\section{Evaluaci\'on}

\subsection{Resultados obtenidos en validaci\'on y test}

La versi\'on actual del MVP soporta seis intenciones anal\'iticas: crecimiento de precio, comparaci\'on de activos, visi\'on descriptiva general, an\'alisis de retornos, riesgo hist\'orico y an\'alisis t\'ecnico. La bater\'ia principal de pruebas incluye ejecuciones correctas y escenarios negativos. En el estado actual del repositorio se documenta una validaci\'on de 9 pruebas superadas sobre 9.

\begin{table}[h]
\centering
\begin{tabular}{llll}
\toprule
\textbf{Tipo} & \textbf{Cobertura} & \textbf{Criterio comprobado} & \textbf{Estado} \\
\midrule
Funcional & 6 intenciones & estado completado y c\'odigo 0 & 6/6 \\
Errores & intenci\'on no soportada & mensaje controlado & 1/1 \\
Errores & CSV inexistente & mensaje controlado & 1/1 \\
Errores & cardinalidad t\'ecnica inv\'alida & mensaje controlado & 1/1 \\
\bottomrule
\end{tabular}
\caption{Resumen de pruebas automatizadas del MVP.}
\end{table}

La validaci\'on funcional cubre los casos \texttt{price\_growth}, \texttt{compare\_assets}, \texttt{asset\_overview}, \texttt{return\_analysis}, \texttt{historical\_risk\_analysis} y \texttt{technical\_analysis}. En todos ellos se comprueba que el workflow termina en estado \texttt{completed}, que el script generado devuelve c\'odigo de salida 0 y que la respuesta contiene informaci\'on esperada del activo o del tipo de an\'alisis.

Tambi\'en se ha probado la integraci\'on del perfil universitario con activaci\'on del LLM solo para interpretaci\'on. En los seis casos representativos el flujo finaliz\'o correctamente, sin fallback, y sin generar recomendaciones directas de inversi\'on. La latencia media observada fue de aproximadamente 6,76 segundos por consulta, con un rango entre 4,71 y 9,87 segundos.

\begin{table}[h]
\centering
\begin{tabular}{lrrrr}
\toprule
\textbf{Perfil} & \textbf{Casos} & \textbf{Completados} & \textbf{Fallbacks} & \textbf{Latencia media} \\
\midrule
Determinista & 6 & 6 & 0 & baja, local \\
Universidad vLLM & 6 & 6 & 0 & 6,76 s \\
\bottomrule
\end{tabular}
\caption{Comparaci\'on resumida entre baseline determinista e interpretaci\'on con LLM.}
\end{table}

\subsection{Curvas de aprendizaje y sobreajuste}

No hay curvas de aprendizaje en el sentido cl\'asico porque no se entrena un modelo predictivo. Sin embargo, s\'i puede analizarse el sobreajuste de prompts o reglas: un pipeline podr\'ia funcionar bien con los ejemplos de desarrollo y fallar con consultas nuevas. Para evitarlo, se propone ampliar el conjunto de test con variaciones de lenguaje, tickers, periodos, intervalos y errores deliberados.

\subsection{Comparaci\'on con m\'etodos base}

La comparaci\'on principal se realiza contra el modo determinista. Este modo act\'ua como l\'inea base porque ofrece respuestas reproducibles y controladas. El LLM se considera superior solo si mejora la claridad, cobertura o adaptabilidad sin introducir errores, recomendaciones financieras ni contradicciones con los resultados num\'ericos. En la validaci\'on actual, el LLM no cambia los c\'alculos: recibe los resultados estructurados ya ejecutados y redacta una explicaci\'on final, lo que reduce el riesgo de alucinaci\'on num\'erica.

\subsection{Pruebas en entorno real}

Para probar el sistema en un entorno real se deber\'ia desplegar una API interna y registrar consultas reales de usuarios acad\'emicos o analistas. Las pruebas incluir\'ian validaci\'on de entradas, tiempo de respuesta, tasa de errores, calidad percibida de las explicaciones y revisi\'on humana de casos con potencial ambig\"uedad. Tambi\'en se deben simular fallos de proveedor LLM, CSV corruptos, latencias elevadas y datos incompletos.

\subsection{Criterios para lanzar a producci\'on}

El rendimiento ser\'ia v\'alido para un entorno real si se cumplen estas condiciones:

\begin{itemize}
    \item cobertura suficiente de intenciones frecuentes;
    \item tasa baja de errores no controlados;
    \item trazabilidad completa de datos, plan, c\'odigo y resultado;
    \item ausencia de recomendaciones de inversi\'on;
    \item latencia aceptable para el usuario;
    \item fallback fiable cuando falle el LLM;
    \item revisi\'on de seguridad para la ejecuci\'on de c\'odigo generado.
\end{itemize}

Las principales limitaciones actuales son el tama\~no reducido de la bater\'ia de evaluaci\'on, la ausencia de pruebas con usuarios reales, la dependencia de proveedores LLM externos cuando se activa la interpretaci\'on generativa y la necesidad de reforzar el entorno de ejecuci\'on antes de permitir generaci\'on libre de c\'odigo. Estas limitaciones no invalidan el MVP, pero s\'i marcan qu\'e faltar\'ia antes de considerarlo un sistema productivo.

\section{Aspectos de ciberseguridad y consideraciones \'eticas}

\subsection{Posibles ataques al modelo}

El sistema puede verse afectado por ataques propios de aplicaciones con LLM:

\begin{itemize}
    \item \textbf{Prompt injection}: el usuario intenta forzar al modelo a ignorar instrucciones o producir asesoramiento financiero.
    \item \textbf{Evasion adversarial}: entradas dise\~nadas para saltarse validaciones o activar un agente incorrecto.
    \item \textbf{Data poisoning}: incorporaci\'on de CSV manipulados para sesgar los resultados.
    \item \textbf{Model stealing}: abuso de la API para inferir comportamiento del modelo o prompts internos.
    \item \textbf{Code injection}: intento de introducir rutas, comandos o fragmentos que acaben en el c\'odigo generado.
\end{itemize}

\subsection{Medidas de protecci\'on}

Las medidas previstas son:

\begin{itemize}
    \item validaci\'on estricta del esquema de entrada;
    \item listas de intenciones soportadas y restricciones por agente;
    \item ejecuci\'on de c\'odigo en proceso controlado y con rutas acotadas;
    \item separaci\'on entre resultados calculados y redacci\'on del LLM;
    \item limitaci\'on de tokens, timeouts y reintentos;
    \item logging de errores y avisos;
    \item fallback determinista ante respuestas no v\'alidas del modelo.
\end{itemize}

\subsection{Privacidad de los datos}

Los datos financieros de mercado no son datos personales. Aun as\'i, si el sistema evolucionara para aceptar carteras, preferencias o historiales de usuario, ser\'ia necesario aplicar minimizaci\'on de datos, anonimizar identificadores y evitar enviar informaci\'on sensible a proveedores externos sin consentimiento.

\subsection{Sesgos y equidad}

Los datos hist\'oricos pueden reflejar sesgos de mercado, periodos an\'omalos o activos con liquidez desigual. El sistema debe explicar que los resultados dependen del periodo analizado y no extrapolar comportamientos pasados como garant\'ia futura. Tambi\'en debe evitar lenguaje categ\'orico que pueda interpretarse como recomendaci\'on personalizada.

\subsection{Uso responsable}

El proyecto se limita al an\'alisis descriptivo. No genera instrucciones de compra o venta, no estima rentabilidad futura y no sustituye asesoramiento financiero profesional. Esta frontera debe aparecer tanto en la documentaci\'on como en las respuestas generadas.

\section{Despliegue}

\subsection{Plataforma de despliegue}

La opci\'on recomendada es un despliegue cloud ligero o un servidor universitario. El backend se expondr\'ia mediante una API REST o una interfaz de l\'inea de comandos evolucionada. Para uso docente, bastar\'ia un servidor local con CPU y acceso configurado a las APIs LLM. Para uso m\'as intensivo, se podr\'ia desplegar en AWS, GCP o Azure.

\subsection{Formato del modelo final}

No existe un \'unico fichero de modelo entrenado. El ``modelo final'' del sistema es la combinaci\'on de c\'odigo, configuraci\'on, agentes, prompts y proveedor LLM seleccionado. Si se incorporara un modelo abierto local, se servir\'ia mediante vLLM y pesos en formato compatible con Hugging Face.

\subsection{Contenedores y orquestaci\'on}

El despliegue se puede empaquetar con Docker, incluyendo dependencias Python, estructura de carpetas, variables de entorno y comandos de ejecuci\'on. Kubernetes solo ser\'ia necesario si se requiere escalado, alta disponibilidad o separaci\'on entre servicios de API, ejecuci\'on y monitorizaci\'on.

\subsection{Monitorizaci\'on en producci\'on}

La monitorizaci\'on debe incluir:

\begin{itemize}
    \item n\'umero de consultas por intenci\'on;
    \item latencia por fase;
    \item tasa de errores y fallbacks;
    \item coste aproximado de llamadas LLM;
    \item proveedor y modelo utilizados;
    \item avisos de calidad de datos;
    \item trazas de ejecuci\'on para auditor\'ia.
\end{itemize}

\section{Costes y sostenibilidad}

\subsection{Costes de entrenamiento y despliegue}

El coste directo para el desarrollo realizado hasta ahora ha sido nulo. Una de las decisiones del MVP ha sido comprobar hasta d\'onde se puede llegar sin gastar dinero: modo determinista en CPU local, datos descargados de fuentes gratuitas, librer\'ias abiertas y uso opcional de APIs solo cuando existe acceso disponible. Adem\'as, no se ha entrenado un modelo desde cero ni se han ajustado pesos, as\'i que no hay coste propio de entrenamiento.

Los costes potenciales aparecer\'ian al escalar el sistema o al usar de forma continuada modelos grandes. En particular, la universidad ha habilitado un nodo donde se aloja un modelo de mayor tama\~no mediante vLLM. Para este trabajo ese recurso no supone un coste directo para el estudiante, pero s\'i tiene un coste real de infraestructura que convendr\'ia solicitar a la universidad antes de la memoria final: tipo de GPU, horas aproximadas de uso, consumo energ\'etico, coste de amortizaci\'on o tarifa equivalente en cloud. Esa estimaci\'on permitir\'ia completar el an\'alisis econ\'omico con mayor precisi\'on.

Los costes principales a considerar en una versi\'on ampliada son:

\begin{itemize}
    \item llamadas a APIs LLM externas;
    \item uso de GPU si se sirve un modelo abierto con vLLM;
    \item almacenamiento de logs y artefactos;
    \item mantenimiento del entorno de ejecuci\'on.
\end{itemize}

Para un uso acad\'emico con pocos casos, el modo determinista y las llamadas ocasionales a API son suficientes. Para un piloto con usuarios reales, conviene estimar coste por consulta y aplicar l\'imites de uso.

\subsection{Uso eficiente de recursos}

La sostenibilidad se mejora evitando llamadas LLM innecesarias. El sistema puede ejecutar c\'alculos en CPU y reservar el modelo de lenguaje para fases donde aporte valor real. Tambi\'en se puede cachear resultados, reutilizar datos descargados, limitar el tama\~no de contexto enviado al LLM y registrar qu\'e tareas se resuelven correctamente sin componente generativo.

\section{Conclusiones}

La idea principal del TFM es construir un sistema multiagente analista que funcione correctamente de extremo a extremo. La contribuci\'on central no es predecir mercados en esta fase, sino dise\~nar e implementar una arquitectura capaz de recibir una consulta estructurada, seleccionar el agente adecuado, cargar datos financieros, generar un plan de an\'alisis, ejecutar c\'alculos verificables y producir una explicaci\'on final. En este sentido, el proyecto integra modelos de Deep Learning dentro de una arquitectura software realista, manteniendo una posici\'on prudente ante el dominio financiero.

El MVP ya cuenta con una base funcional: seis intenciones anal\'iticas, pruebas automatizadas, gesti\'on de errores, generaci\'on y ejecuci\'on de scripts, y una capa LLM configurable. Los siguientes pasos naturales son ampliar el conjunto de evaluaci\'on, reforzar la seguridad de la ejecuci\'on de c\'odigo, registrar experimentos con una herramienta MLOps y comparar de forma sistem\'atica el modo determinista frente a distintos proveedores LLM.

Como trabajo futuro inmediato, el sistema deber\'ia incorporar una capa de procesamiento de lenguaje natural que permita recibir peticiones en texto libre. Esa capa tendr\'ia que extraer intenci\'on, tickers, periodos, intervalos y restricciones de la consulta del usuario, para convertirlas en la entrada estructurada que el MVP ya sabe ejecutar. Tambi\'en se plantea automatizar la obtenci\'on de datos mediante \texttt{yfinance}, de modo que el sistema pueda descargar por s\'i mismo las series necesarias cuando no existan en local, registrando proveedor, fecha de descarga y metadatos para conservar la reproducibilidad.

A m\'as largo plazo, una vez consolidado el sistema analista y su trazabilidad, se podr\'ia estudiar la incorporaci\'on de modelos de \textit{forecasting}. Esta extensi\'on deber\'ia tratarse como una fase diferenciada, con divisi\'on temporal de datos, m\'etricas predictivas, comparaci\'on con baselines cl\'asicos y comunicaci\'on clara de incertidumbre, evitando que el sistema descriptivo actual se confunda con asesoramiento financiero o predicci\'on garantizada.

\begin{thebibliography}{9}

\bibitem{react}
Yao, S. et al. \textit{ReAct: Synergizing Reasoning and Acting in Language Models}. arXiv:2210.03629, 2022.
\url{https://arxiv.org/abs/2210.03629}

\bibitem{toolformer}
Schick, T. et al. \textit{Toolformer: Language Models Can Teach Themselves to Use Tools}. arXiv:2302.04761, 2023.
\url{https://arxiv.org/abs/2302.04761}

\bibitem{autogen}
Wu, Q. et al. \textit{AutoGen: Enabling Next-Gen LLM Applications via Multi-Agent Conversation}. arXiv:2308.08155, 2023.
\url{https://arxiv.org/abs/2308.08155}

\bibitem{fingpt}
Yang, H. et al. \textit{FinGPT: Open-Source Financial Large Language Models}. arXiv:2306.06031, 2023.
\url{https://arxiv.org/abs/2306.06031}

\bibitem{mlflow}
Zaharia, M. et al. \textit{Accelerating the Machine Learning Lifecycle with MLflow}. IEEE Data Engineering Bulletin, 2018.
\url{https://mlflow.org/}

\bibitem{yfinance}
Aroussi, R. \textit{yfinance Python Library}. GitHub repository.
\url{https://github.com/ranaroussi/yfinance}

\end{thebibliography}

\end{document}
